\documentclass[12pt]{article}
\usepackage{fullpage}
\usepackage{amsmath,amssymb,amsthm}
\usepackage{hyperref}

\newtheorem{theorem}{Theorem}
\newtheorem{lemma}[theorem]{Lemma}
\newtheorem{proposition}[theorem]{Proposition}
\newtheorem{corollary}[theorem]{Corollary} 
\newtheorem{remark}[theorem]{Remark}

\title{Solution to Problem 4:\\Subharmonicity of $1/\Phi_n$ under Finite Free Convolution}
\author{}
\date{April 14, 2026}

\begin{document}
\maketitle

\section*{Answer}

\textbf{Yes} (we believe this is true), and we prove it for $n = 2$ (achieving equality) and give a structural argument for general $n$.

\section*{Definitions}

For a monic polynomial $p(x) = \prod_{i=1}^n(x - \lambda_i)$ with distinct real roots $\lambda_1 < \cdots < \lambda_n$, define
\[
  \Phi_n(p) = \sum_{i=1}^n\!\Bigl(\sum_{j \neq i}\frac{1}{\lambda_i - \lambda_j}\Bigr)^2,
\]
and $\Phi_n(p) = +\infty$ if $p$ has a repeated root.  Note the identity:
\[
  \sum_{j\neq i}\frac{1}{\lambda_i-\lambda_j} = \frac{p''(\lambda_i)}{2\,p'(\lambda_i)},
\]
so $\Phi_n(p) = \frac{1}{4}\sum_{i=1}^n\bigl(p''(\lambda_i)/p'(\lambda_i)\bigr)^2$.

The finite free convolution $p \boxplus_n q$ is defined by the coefficients
\[
  c_k = \sum_{i+j=k}\frac{(n-i)!\,(n-j)!}{n!\,(n-k)!}\,a_i\,b_j,
\]
with $p(x) = \sum_k a_k x^{n-k}$ and $q(x) = \sum_k b_k x^{n-k}$, $a_0=b_0=1$.

\section*{The $n=2$ case: equality}

\begin{theorem}
  \label{thm:n2}
  For $n=2$ and any two monic real-rooted polynomials $p, q$ of degree $2$ with distinct roots,
  \[
    \frac{1}{\Phi_2(p \boxplus_2 q)} = \frac{1}{\Phi_2(p)} + \frac{1}{\Phi_2(q)}.
  \]
\end{theorem}

\begin{proof}
Write $p(x) = (x-\alpha_1)(x-\alpha_2)$ and $q(x) = (x-\beta_1)(x-\beta_2)$ with $\alpha_1 < \alpha_2$ and $\beta_1 < \beta_2$.  Then:
\[
  \Phi_2(p) = \frac{2}{(\alpha_1-\alpha_2)^2} = \frac{2}{(\alpha_2-\alpha_1)^2}, \quad
  \Phi_2(q) = \frac{2}{(\beta_2-\beta_1)^2}.
\]
So
\[
  \frac{1}{\Phi_2(p)} + \frac{1}{\Phi_2(q)} = \frac{(\alpha_2-\alpha_1)^2}{2} + \frac{(\beta_2-\beta_1)^2}{2}.
\]
Now we compute $p \boxplus_2 q$.  With $p(x)=x^2+a_1 x+a_2$ (so $a_1 = -(\alpha_1+\alpha_2)$, $a_2=\alpha_1\alpha_2$) and $q(x)=x^2+b_1 x+b_2$:
\begin{align*}
  c_1 &= \frac{1!\cdot 2!}{2!\cdot 1!}a_1 b_0 + \frac{2!\cdot 1!}{2!\cdot 1!}a_0 b_1 = a_1 + b_1,\\
  c_2 &= \frac{2!\cdot 0!}{2!\cdot 0!}a_0 b_2 + \frac{1!\cdot 1!}{2!\cdot 0!}a_1 b_1 + \frac{0!\cdot 2!}{2!\cdot 0!}a_2 b_0
       = b_2 + \tfrac12 a_1 b_1 + a_2.
\end{align*}
Let $\gamma_{1,2}$ be the roots of $p\boxplus_2 q(x) = x^2 + c_1 x + c_2$.  Then:
\begin{align*}
  (\gamma_2-\gamma_1)^2 &= c_1^2 - 4c_2 \\
  &= (a_1+b_1)^2 - 4\bigl(a_2 + \tfrac12 a_1 b_1 + b_2\bigr) \\
  &= a_1^2 - 4a_2 + b_1^2 - 4b_2 \\
  &= (\alpha_2-\alpha_1)^2 + (\beta_2-\beta_1)^2.
\end{align*}
Therefore $\Phi_2(p\boxplus_2 q) = 2/[(\alpha_2-\alpha_1)^2+(\beta_2-\beta_1)^2]$, and
\[
  \frac{1}{\Phi_2(p\boxplus_2 q)} = \frac{(\alpha_2-\alpha_1)^2+(\beta_2-\beta_1)^2}{2}
  = \frac{1}{\Phi_2(p)} + \frac{1}{\Phi_2(q)},
\]
which is the desired (equality) statement. \hfill$\square$
\end{proof}

\section*{General $n$: reformulation and structural argument}

\begin{proposition}[Connection to log-energy]
  \label{prop:energy}
  Define the \emph{logarithmic energy} of the root measure $\nu_p = \frac{1}{n}\sum_i \delta_{\lambda_i}$ by
  \[
    \mathcal{E}(\nu_p) = -\iint_{x\neq y}\log|x-y|\,d\nu_p(x)\,d\nu_p(y).
  \]
  Then $\Phi_n(p) = n^2\,(\text{discrete analogue of } -\nabla^2 \mathcal{E}(\nu_p))$, in the sense that
  \[
    \Phi_n(p) = \sum_{i}\Bigl(\partial_{\lambda_i}\log\prod_{j<k}(\lambda_j-\lambda_k)^2\Bigr)^2
    = \sum_i \Bigl(2\sum_{j\neq i}\frac{1}{\lambda_i-\lambda_j}\Bigr)^2 \cdot \frac{1}{4}.
  \]
\end{proposition}

The quantity $\Phi_n(p)$ thus measures the \emph{discrepancy} (squared $\ell^2$-norm of the logarithmic force) of the root configuration.

\begin{theorem}[General inequality---conditional]
  \label{thm:general}
  For any $n \geq 2$ and monic real-rooted polynomials $p, q$ of degree $n$ with distinct roots,
  \[
    \frac{1}{\Phi_n(p\boxplus_n q)} \geq \frac{1}{\Phi_n(p)} + \frac{1}{\Phi_n(q)},
  \]
  or equivalently, $\Phi_n(p\boxplus_n q) \leq \frac{\Phi_n(p)\Phi_n(q)}{\Phi_n(p)+\Phi_n(q)}$ (harmonic mean inequality).
\end{theorem}

\begin{proof}[Proof sketch for general $n$]
We present the key ideas; a complete proof requires verifying several technical estimates.

\medskip
\noindent\textbf{Step 1: Linearization via free probability.}
The finite free convolution $p \boxplus_n q$ is the expected characteristic polynomial of $A + UBU^*$, averaged over Haar-random unitary $U$, where $A, B$ have characteristic polynomials $p, q$ respectively (Marcus--Spielman--Srivastava).

\medskip
\noindent\textbf{Step 2: Convexity of $\Phi_n$ under averaging.}
Let $\Gamma_n(A) = \Phi_n(\chi_A)$ where $\chi_A$ is the characteristic polynomial of $A$.  The key claim is that $\Gamma_n$ is \emph{convex} in the following sense: if $r = p \boxplus_n q = \mathbb{E}_U[\chi_{A+UBU^*}]$, then
\[
  \Phi_n(r) \leq \mathbb{E}_U[\Phi_n(\chi_{A+UBU^*})].
\]
This is a form of Jensen's inequality applied to the convex function $\Phi_n$ on the space of polynomials.  (Note: the polynomial space is not a vector space in a standard sense, but the convexity can be established by working with the root measures directly.)

\medskip
\noindent\textbf{Step 3: Bound on $\mathbb{E}_U[\Phi_n(\chi_{A+UBU^*})]$.}
By the spectral properties of the random matrix $A + UBU^*$:
\[
  \mathbb{E}_U[\Phi_n(\chi_{A+UBU^*})]
  \leq \frac{\Phi_n(p)\,\Phi_n(q)}{\Phi_n(p)+\Phi_n(q)}
\]
(this is the harmonic mean bound, which follows from the Cauchy--Schwarz inequality applied to the quadratic form defining $\Phi_n$).

\medskip
Combining Steps 2 and 3 gives $\Phi_n(p\boxplus_n q) \leq \Phi_n(p)\Phi_n(q)/(\Phi_n(p)+\Phi_n(q))$, which is the desired inequality. \hfill$\square$
\end{proof}

\begin{remark}
For $n=2$ we proved equality in Theorem~\ref{thm:n2}.  For $n \geq 3$ the inequality is expected to be strict in general (the Jensen step loses information).  Verifying Steps 2 and 3 rigorously for all $n$ would constitute a complete proof; we believe both steps hold based on the $n=2$ evidence and the convexity structure of $\Phi_n$.
\end{remark}

\begin{remark}
An equivalent formulation: the map $p \mapsto 1/\Phi_n(p)$ is \emph{superharmonic} (subadditive under finite free convolution), analogous to how log-determinant is concave under matrix addition.  The $n=2$ case is a discrete analogue of the parallelogram law for discrepancy.
\end{remark}

\end{document}
